\chapter{The Axioms of Set Theory} \phantomsection \label{ch:zfc-axioms}

The following is a list of axioms from which every other statement in set theory can be deduced.

\begin{axm}[Axiom of Existence]
    \label{axm:axiom-of-existence}
    There exists a set without any elements.
    \begin{equation*}
        \exists A \, \forall x \, \mleft( x \notin A \mright)
    \end{equation*}
\end{axm}

\begin{axm}[Axiom of Regularity]
    \label{axm:axiom-of-regularity}
    If a set \( A \) has at least one element, then there exists an element \( x \) in \( A \) such that
    \( A \) and \( x \) have no element in common.
    \begin{equation*}
        \forall A \, \Big[ \exists x \, \mleft( x \in A \mright)
        \implies
        \exists x \, \big( x \in A \land \forall y \, \mleft( y \in x \implies y \notin A \mright) \big)
        \Big]
    \end{equation*}
\end{axm}

\begin{axm}[Axiom of Extensionality]
    \label{axm:axiom-of-extensionality}
    If every element of \( A \) is an element of \( B \) and every element of \( B \) is also an element of \( A \), then \( A = B \).
    \begin{equation*}
        \forall A,B \, \Big[\, \forall x \left( x \in A \iff x \in B \right) \implies \mleft( A = B \mright) \,\Big]
    \end{equation*}
\end{axm}

\begin{axm}[Axiom of Pairing]
    \label{axm:axiom-of-pairing}
    For any two sets \( a \) and \( b \), there exists a set \( S \) such that \( S \) contains exactly \( a \) and \( b \).
    \begin{equation*}
        \forall a,b \, \, \exists S \, \forall x \, \Big[ x \in S \iff (x = a \,\,{\lor}\,\, x = b) \Big]
    \end{equation*}
\end{axm}

\begin{axm}[Axiom Schema of Specification]
    \label{axm:axiom-schema-of-specification}
    If \( \varphi \) is a well-formed formula of any first-order language of set theory with \( v_{0},\dots, v_{n}\), \( x \), and \( A \)
    being the only \aref{dfn:first-order-language-recursive-definition-of-free-variables}{free variables of \( \varphi \)}, then the following sentence is true:
    \begin{equation*}
        \forall v_{0},\dots,v_{n} \, \forall A \, \exists S \, \forall x \, \Big[ x \in S \iff x \in A \,\,{\land}\,\, \varphi \Big]
    \end{equation*}
\end{axm}

\begin{axm}[Axiom of Union]
    \label{axm:axiom-of-union}
    For any set \( A \), there exists a set \( U \) such that \( x \in U \) if and only if \( x \in S \) for some \( S \in A \).
    \begin{equation*}
        \forall A \, \exists U \, \forall x \, \Big[ x \in U \iff \exists S \, \big( S \in A \,\,{\land}\,\, x \in S \big) \Big]
    \end{equation*}
\end{axm}

\begin{axm}[Axiom of Power Set]
    \label{axm:axiom-of-power-set}
    For any set \( A \), there exists a set \( P \) such that for any \( x \), \( x \in P \) if and only if \,\( \forall y \, \bigl( y \in x \implies y \in A \bigr) \)
    \begin{equation*}
        \forall A \, \exists P \, \forall x \, \Big[ x \in P \iff \forall y \, \big( y \in x \implies y \in A \big)\Big]
    \end{equation*}
\end{axm}

\begin{axm}[Axiom of Infinity]
    \label{axm:axiom-of-infinity}
    An inductive set exists.
    \begin{equation*}
        \exists I \, \Big[ \emptyset \in I \land \forall x \, \big( x \in I \implies x \union \{x\} \in I\big)\Big]
    \end{equation*}
\end{axm}
